\documentclass[conference]{IEEEtran}

\usepackage[T1]{fontenc}
\usepackage[utf8]{inputenc}
\usepackage[brazil]{babel}
\usepackage{csquotes}
\usepackage{booktabs}
\usepackage{array}
\usepackage{graphicx}
\usepackage{siunitx}
\usepackage[hidelinks]{hyperref}

\title{Economia de Combustíveis Fósseis no Brasil e no Mundo: Uma Análise Comparativa de Petróleo, Gás Natural e Carvão Mineral}

\author{\IEEEauthorblockN{Matheus Faria}
\IEEEauthorblockA{Curso de Engenharia Elétrica \\
Universidade Federal de Juiz de Fora \\
Juiz de Fora, Brasil \\
\texttt{matheus@example.com}}
\and
\IEEEauthorblockN{Grupo 7 -- Economia da Energia}
\IEEEauthorblockA{Disciplina de Economia da Energia \\
Universidade Federal de Juiz de Fora \\
Juiz de Fora, Brasil}}

\begin{document}

\maketitle

\begin{abstract}
Este artigo apresenta a estrutura inicial de um estudo comparativo sobre a economia de petróleo, gás natural e carvão mineral, articulando o contexto do Brasil com o cenário mundial. São definidos o problema de pesquisa, os objetivos, o recorte analítico e um protocolo de indicadores comuns para garantir comparabilidade entre fontes, períodos e unidades.
\end{abstract}

\begin{IEEEkeywords}
Economia da energia, combustíveis fósseis, petróleo, gás natural, carvão mineral, Brasil, mercado internacional.
\end{IEEEkeywords}

\section{Introdu\c{c}\~ao}

A economia dos combust\'iveis f\'osseis permanece central para o funcionamento dos sistemas energ\'eticos e para a competitividade industrial. Este artigo investiga, de forma comparativa, petr\'oleo, g\'as natural e carv\~ao mineral, combinando evid\^encias para o Brasil e para o mundo com base em fontes oficiais e setoriais \cite{ben2026,ei2026}.

\section{Justificativa}

A compara\c{c}\~ao integrada entre os tr\^es combust\'iveis permite identificar diferen\c{c}as de estrutura de mercado, forma\c{c}\~ao de pre\c{c}os e padr\~oes de seguran\c{c}a energ\'etica. O estudo tamb\'em contribui para discutir a transi\c{c}\~ao energ\'etica sem perder de vista restri\c{c}\~oes econ\^omicas, institucionais e geopol\'iticas.

\section{Tema, recorte e objetivos}

\subsection{Tema e t\'itulo provis\'orio}
Tema: economia comparada dos combust\'iveis f\'osseis. T\'itulo provis\'orio: \emph{Economia de Combust\'iveis F\'osseis no Brasil e no Mundo: Uma An\'alise Comparativa de Petr\'oleo, G\'as Natural e Carv\~ao Mineral}.

\subsection{Pergunta de pesquisa}
Como diferem, entre Brasil e mundo, a din\^amica de reservas, produ\c{c}\~ao, consumo, com\'ercio e pre\c{c}os de petr\'oleo, g\'as natural e carv\~ao mineral, e quais implica\c{c}\~oes econ\^omicas decorrem dessas diferen\c{c}as?

\subsection{Objetivo geral}
Comparar a economia de petr\'oleo, g\'as natural e carv\~ao mineral no Brasil e no mundo com base em indicadores padronizados e compar\'aveis.

\subsection{Objetivos espec\'ificos}
\begin{itemize}
  \item consolidar um conjunto comum de indicadores para os tr\^es combust\'iveis;
  \item analisar diferen\c{c}as de oferta, demanda e com\'ercio entre Brasil e mundo;
  \item comparar a forma\c{c}\~ao de pre\c{c}os e fatores de volatilidade;
  \item discutir implica\c{c}\~oes para seguran\c{c}a energ\'etica e planejamento setorial.
\end{itemize}

\subsection{Recorte}
\textbf{Combust\'iveis}: petr\'oleo, g\'as natural e carv\~ao mineral. \\
\textbf{Abrang\^encia geogr\'afica}: Brasil e mundo. \\
\textbf{Horizonte temporal}: s\'eries hist\'oricas recentes (prioridade para 2015--2025, conforme disponibilidade) e proje\c{c}\~oes quando explicitamente indicadas pelas fontes.

\subsection{Or\c{c}amento de conte\'udo (meta: 17--22 p\'aginas)}
\begin{table}[h]
\centering
\caption{Distribui\c{c}\~ao planejada de p\'aginas}
\begin{tabular}{lc}
\toprule
Se\c{c}\~ao & Faixa (p\'aginas) \\
\midrule
Introdu\c{c}\~ao + justificativa & 2--3 \\
Tema, recorte e m\'etodo & 2--3 \\
An\'alise do petr\'oleo & 3--4 \\
An\'alise do g\'as natural & 3--4 \\
An\'alise do carv\~ao mineral & 3--4 \\
Discuss\~ao comparativa & 2--3 \\
Conclus\~ao & 1--2 \\
\midrule
Total previsto & 17--22 \\
\bottomrule
\end{tabular}
\end{table}

\section{M\'etodo e indicadores comuns}

O artigo adota estrat\'egia descritiva-comparativa, com triangula\c{c}\~ao de fontes oficiais nacionais e internacionais. Para comparabilidade, cada indicador deve registrar fonte, per\'iodo, unidade e abrang\^encia.

\begin{table*}[t]
\centering
\caption{Indicadores comuns para petr\'oleo, g\'as natural e carv\~ao mineral}
\begin{tabular}{p{3.2cm}p{4.0cm}p{2.5cm}p{5.4cm}}
\toprule
Indicador & Unidade de refer\^encia & Periodicidade & Crit\'erio de comparabilidade \\
\midrule
Reservas provadas & barris, bcm ou Mt & anual & convers\~ao para equival\^encia energ\'etica quando necess\'ario \\
Produ\c{c}\~ao & bbl/dia, bcm/ano ou Mt/ano & anual/mensal & explicitar m\'etrica f\'isica e, se usado, equivalente energ\'etico \\
Consumo (demanda) & bbl/dia, bcm/ano ou Mt/ano & anual/mensal & manter distin\c{c}\~ao entre consumo final e demanda prim\'aria \\
Com\'ercio internacional & exporta\c{c}\~oes/importa\c{c}\~oes (volume e valor) & anual & separar balan\c{c}a f\'isica da monet\'aria \\
Pre\c{c}os de refer\^encia & US\$/bbl, US\$/MMBtu, US\$/t & mensal/anual & informar pre\c{c}o nominal, moeda e deflator quando aplic\'avel \\
\bottomrule
\end{tabular}
\end{table*}

Per\'iodo-base comum priorizado: 2015--2025 para s\'eries hist\'oricas, com extens\~ao condicionada \`a disponibilidade e consist\^encia das fontes.

\section{An\'alise setorial: petr\'oleo}

Se\c{c}\~ao reservada para integrar dados de reservas, produ\c{c}\~ao, consumo, com\'ercio e pre\c{c}os de petr\'oleo no Brasil e no mundo, seguindo o protocolo de indicadores comuns.

\section{An\'alise setorial: g\'as natural}

Se\c{c}\~ao reservada para integrar dados de reservas, produ\c{c}\~ao, consumo, com\'ercio e pre\c{c}os de g\'as natural no Brasil e no mundo, seguindo o protocolo de indicadores comuns.

\section{An\'alise setorial: carv\~ao mineral}

Se\c{c}\~ao reservada para integrar dados de reservas, produ\c{c}\~ao, consumo, com\'ercio e pre\c{c}os de carv\~ao mineral no Brasil e no mundo, seguindo o protocolo de indicadores comuns.

\section{Discuss\~ao}

Se\c{c}\~ao pronta para compara\c{c}\~ao transversal dos tr\^es mercados, destacando converg\^encias, diverg\^encias e implica\c{c}\~oes econ\^omicas para o planejamento energ\'etico.

\section{Conclus\~ao}

Se\c{c}\~ao pronta para s\'intese dos resultados, limita\c{c}\~oes do estudo e encaminhamentos para a vers\~ao final do artigo e da apresenta\c{c}\~ao.


\bibliographystyle{IEEEtran}
\bibliography{referencias}

\end{document}
